\section*{الفصل \ref{ch:b3:structural-biology} --- البيولوجيا البنيوية للبروتينات}

\begin{solution}{exo:b3:structural-biology:1}
الأولية: تسلسل البقايا البالغة $141$ في سلسلة $\alpha$.
والثانوية: لوالبها $\alpha$ الثمانية. والثالثية: \omterm{def:b3:structural-biology:levels}{طية} الغلوبين في سلسلة
واحدة ملتفة حول هيمها. والرابعية: الرباعي $\alpha_{2}\beta_{2}$
وطريقة دوران ثنائييه بعضهما على بعض.
\end{solution}

\begin{solution}{exo:b3:structural-biology:2}
$36\times \qty{0.15}{nm} = \qty{5.4}{nm}$؛ و$36/3.6 = 10$ لفات؛
و$36 - 4 = 32$ رابطة هيدروجينية.
\end{solution}

\begin{solution}{exo:b3:structural-biology:3}
أن البنية الأصلية، بما فيها الازدواج الصحيح لأربع
روابط كبريتية من بين $105$ إمكانًا، تُستعاد عفويًا من
السلسلة غير المطوية: فالتسلسل يحوي كل المعلومة اللازمة،
والحالة الأصلية هي الأصغرية الترموديناميكية. وإجراؤها في أنبوب
اختبار ببروتين نقي استبعد أي قالب أو إنزيم أو آلة
خلوية مصدرًا للمعلومة.
\end{solution}

\begin{solution}{exo:b3:structural-biology:4}
الرنين المغناطيسي النووي: البروتينات الصغيرة، دون نحو \qty{40}{kDa}، في المحلول.
\omterm{def:b3:structural-biology:methods}{وحيود الأشعة السينية}: أي حجم يتبلور، من الببتيدات إلى
الريبوزومات. \omterm{def:b3:structural-biology:methods}{والمجهر الإلكتروني المبرَّد}: المعقّدات الكبيرة، فوق نحو
\qty{50}{kDa}، بلا بلورات. والرنين المغناطيسي يبلّغ عن الحركة مباشرةً (و
المجهر المبرَّد يستطيع فرز الجزيئات في أصناف انتظامية)؛ أما البنية البلورية
فمتوسط ساكن.
\end{solution}

\begin{solution}{exo:b3:structural-biology:5}
$2^{150} \approx 1.4\times 10^{45}$ انتظام، أي $1.4\times 10^{33}$
ثانية $\approx 5\times 10^{25}$ سنة. و$3^{20} \approx 3.5\times 10^{9}$،
أي $3.5$ مليثانية: فالببتيد يستطيع البحث استقصاءً في ثانية؛ أما
البروتين فلا يستطيعه في عمر الكون.
\end{solution}

\begin{solution}{exo:b3:structural-biology:6}
$Y = [\alpha(1+\alpha) + Lc\alpha(1+c\alpha)]/[(1+\alpha)^{2} +
L(1+c\alpha)^{2}]$. فعند $\alpha = 1$: $3.01/106 = 0.028$ (والموقع الواحد
$0.50$). وعند $\alpha = 10$: $121/242 = 0.50$ (و$0.91$). وعند $\alpha = 100$:
$\num{10300}/\num{10601} = 0.97$ (و$0.99$). فمنحني MWC يتخلف كثيرًا
عن الموقع الواحد عند الرابطة القليلة، ثم يرتفع بحدة --- من $0.03$
إلى $0.97$ على عقدين يذهب فيهما الموقع الواحد من $0.5$ إلى
$0.99$: وتلك الحدة هي \omterm{def:b3:structural-biology:allostery}{التعاونية}.
\end{solution}

\begin{solution}{exo:b3:structural-biology:7}
لا يرتبط ثنائي فوسفوغليسيرات إلا \omterm{def:b3:structural-biology:allostery}{بالحالة T} (في التجويف المركزي)، فهو
يثبّت T: أي يرفع في النموذج $L$، وهي نسبة $T/R$ في
البروتين بلا رابطة، فيزيح المنحني إلى اليمين --- أي ألفة
أقل، وأكسجين أكثر يُطلق عند ضغوط الأنسجة. والكريات الحمر المخزونة
تفقد المركّب، فتهبط $L$، وينزاح المنحني إلى اليسار، ويمسك الدم
المنقول أكسجينه بدل أن يتخلى عنه في الأنسجة، حتى
تعيد الخلايا تركيبه على مدى يوم.
\end{solution}

\begin{solution}{exo:b3:structural-biology:8}
$d = \lambda/(2\sin\theta) = 0.1/(2\times 0.259) = \qty{0.19}{nm}$: أي
يوضع الهيكل وكل سلسلة جانبية، وتُميَّز الذرات المفردة.
ولبلوغ \qty{0.12}{nm}: $\sin\theta = 0.1/0.24 = 0.417$، ومنه $\theta =
\qty{25}{\degree}$.
\end{solution}

\begin{solution}{exo:b3:structural-biology:9}
الليزين المشحون المدفون بين سلاسل جانبية غير قطبية يكلّف عشرات
الكيلوجولات في المول (فشحنته غير مذابة، ودفن الليوسين الكاره للماء
مفقود)، أي أكثر من $\Delta G$ الطي كله: فينفتل البروتين
أو يعيد اللب ترتيبه. أما على السطح فيكون الليزين
مُماهًا كما كان ليكون في الحالة غير المطوية ولا شيء يُفقد.
و$\qty{-30}{kJ/mol} \approx -12RT$ تعني $K = e^{12} \approx 1.6\times
10^{5}$: أي جزيء واحد من \num{160000} غير مطوي في أي لحظة، و
الهامش كله يساوي رابطتين أو ثلاثًا هيدروجينية من مئات
--- أي قدر طفرة واحدة.
\end{solution}

\begin{solution}{exo:b3:structural-biology:10}
مجموع أنواع $R$ هو $[R_{0}](1+\alpha)^{n}$ ومجموع أنواع $T$ هو
$[R_{0}]L(1+c\alpha)^{n}$، فتكون $\bar R = (1+\alpha)^{n}/[(1+\alpha)^{n}
+ L(1+c\alpha)^{n}]$. ومن أجل $c \to 0$ يصير حد $T$ هو $L$، وتكون $\bar R =
1/2$ حين $(1+\alpha)^{n} = L$، أي $\alpha = L^{1/n} - 1$. وفي
الهيموغلوبين $(3\times 10^{5})^{1/4} = 23.4$، فتكون $\alpha \approx 22$:
أي \qty{22}{mmHg}، وهي قريبة من $P_{50}$ البالغ \qty{26}{mmHg} ---
فالمفتاح الانتظامي ونصف الإشباع يكادان يتطابقان.
\end{solution}

\begin{solution}{exo:b3:structural-biology:11}
قوبلت بالمقاومة لأن كل عامل معدٍ معروف كان يتضاعف
عبر حمض نووي، والبروتين لا يستطيع نسخ تسلسله. و
وجودُ حمض نووي لازم، أو بيانُ عجز بروتين منقّى عن
النقل، كان ليدحضها. وأما الدعم: فالعدوائية تنجو من النوكليازات،
ومن الضوء فوق البنفسجي والإشعاع المؤيّن بجرعات تهدم أي جينوم،
لكنها تُهدم بمُمسِّخات البروتين وبالبروتيازات؛ وPrP
المؤتلف المطوي خارج الجسم ليفاتٍ ينقل المرض إلى الحيوانات، بخصائص
السلالة نفسها عند التمرير. والبذرة لازمة لأن
تحويل الوحيد السوي بطيء بطئًا مانعًا وحده --- إذ يجب أولًا أن
تتكوّن نواة بنية التصالب $\beta$، وهي الحدث المحدد
للسرعة --- أما نهاية ليف متكوّن سلفًا فتحوّل
الوحيدات سريعًا: فالشكل يُنسخ بنمو موجَّه بالقالب.
\end{solution}

\begin{solution}{exo:b3:structural-biology:12}
تسلسلات كثيرة تنطوي إلى البنية نفسها: فما تحتاجه \omterm{def:b3:structural-biology:levels}{الطية}
نمطٌ من المواضع الكارهة للماء والقطبية، وقليل من الغليسينات والبرولينات
عند المنعطفات، والبقايا التي تؤدي العمل؛ وكل ما عدا ذلك على
السطح يمكن أن يتغير بلا جزاء، فتتراكم الطفرات هناك بينما
يحفظ الانتقاء \omterm{def:b3:structural-biology:levels}{الطية} والوظيفة. والبقايا المحفوظة
إذن هي اللب (نمطًا أكثر منها هويات)، و
البقايا الحفّازة والرابطة، والغليسينات والبرولينات والسيستئينات
البنيوية. وطرائق الملامح تزن تلك الأعمدة بالضبط فتتعرف
أفراد الأسرة عند \qty{15}{\%} تطابقًا؛ وتصميم البروتينات يجري
المنطق عكسًا، فيختار نمطًا كارهًا للماء \omterm{def:b3:structural-biology:levels}{لطية} مستهدفة
ويدع السطح يكون أي شيء.
\end{solution}

\begin{solution}{pb:b3:structural-biology:1}
\textbf{1.} $0.75\times 150 \approx 112$ بقية لولبية؛ و$112\times
\qty{0.15}{nm} = \qty{17}{nm}$؛ و$112/3.6 \approx 31$ لفة.
\textbf{2.} $112 - 8\times 4 = 80$ رابطة هيدروجينية.
\textbf{3.} الكتلة $150\times 110 = \qty{16500}{Da} = 2.7\times 10^{-20}$
غرام؛ والحجم $2.7\times 10^{-20}\times 0.73 = 2.0\times 10^{-20}$ سم$^{3}$
$= \qty{20}{nm^3}$؛ ونصف القطر $(3V/4\pi)^{1/3} = \qty{1.7}{nm}$.
\textbf{4.} الممتدة $150\times 0.36 = \qty{54}{nm}$ في مواجهة قطر
قدره \qty{3.4}{nm}: أي عامل $16$.
\textbf{5.} البقايا اللولبية \qty{75}{\%} من الطول
الممتد، أي $112\times 0.36 = \qty{40}{nm}$؛ ويضغطها اللولب إلى
\qty{17}{nm}، أي بعامل $0.36/0.15 = 2.4$.
\textbf{6.} $40\times 4 = \qty{160}{kJ/mol}$ مكسوبة، و$150\times 0.9 =
\qty{135}{kJ/mol}$ مفقودة: فالصافي $\Delta G \approx \qty{-25}{kJ/mol}$، في
المدى الهامشي.
\textbf{7.} $3^{150} \approx 4\times 10^{71}$؛ و$4\times 10^{59}$ ثانية
$\approx 10^{52}$ سنة.
\textbf{8.} $10^{-2}\times 10^{12} = 10^{10}$ انتظام، أي نسبة
$3\times 10^{-62}$ من العدّ.
\textbf{9.} اللوالب في \qty{1}{\micro s} (بالتوازي)؛ و$8! = \num{40320}$
تراصّ عند $10^{6}$ في الثانية تستغرق \qty{40}{ms}: فالمجموع نحو
\qty{40}{ms}، وهي رتبة القدر المرصودة. نعم: فحلّ الأجزاء
استقلالًا ثم الكل قمعٌ --- إذ يُحلَّل فضاء البحث
إلى عوامل، لا يُستقصى.
\textbf{10.} الدورات المتوقعة $1/0.3 = 3.3$، أي نحو \qty{33}{s}؛ وبعد
ست دورات يبقى $0.7^{6} = 0.12$ غير مطوي.
\textbf{11.} $\Delta G = -RT\ln K$ مع $K = 10^{4} \to 10^{2}$: أي من
$-23.7$ إلى \qty{-11.9}{kJ/mol}، أي فقد \qty{12}{kJ/mol} من الثبات
($RT = \qty{2.58}{kJ/mol}$، و$\ln 100 = 4.6$).
\textbf{12.} يبدأ التكتل بنواة، ويرتفع معدل تكوّنها
بأسّ عالٍ من تركيز النوع غير المطوي المهيّأ
للتكتل؛ فمئة ضعف من ذلك النوع ترفع احتمال
تكوّن بذرة خلال عمر برتب من المقدار، ومتى وُجدت
بذرة كان النمو سريعًا ذاتي القولبة.
\textbf{13.} عند $\alpha = 100$: $Y = 1.054\times 10^{8}/1.089\times 10^{8} =
0.97$. وعند $\alpha = 20$: البسط $185\,220 + 103\,680 = 288\,900$،
والمقام $194\,481 + 622\,080 = 816\,561$، ومنه $Y = 0.35$.
\textbf{14.} $\bar R = 1.041\times 10^{8}/1.089\times 10^{8} = 0.96$ في
الرئة؛ و$194\,481/816\,561 = 0.24$ في العضلة.
\textbf{15.} $0.97 - 0.35 = 0.62$ من السعة مفرَّغة (أي ثلثا
الحمولة). وبموقع واحد: $100/126 - 20/46 = 0.79 - 0.43 = 0.36$.
\textbf{16.} السعة $150\times 1.34 = \qty{200}{mL/L}$؛ وإلى العضلة
$0.62\times 200 \approx \qty{125}{mL/L}$. وفي الراحة $(0.97 - 0.78)\times
200 = \qty{38}{mL/L}$؛ و$250/38 \approx \qty{6.5}{L/min}$ من جريان
الدم --- وهي رتبة نتاج القلب في الراحة.
\textbf{17.} عند $L = 3\times 10^{6}$ و$\alpha = 40$: البسط $2.76\times
10^{6} + 3.29\times 10^{6} = 6.05\times 10^{6}$، والمقام $2.83\times
10^{6} + 11.5\times 10^{6} = 14.4\times 10^{6}$، ومنه $Y = 0.42$ بدل
$0.78$ --- أي أكسجين أكثر بكثير يُفرَّغ في الراحة، وهو تكيف الدم مع
الارتفاع والعضلة مع العمل.
\textbf{18.} الجنيني، عند $L = 3\times 10^{4}$ و$\alpha = 30$: البسط
$893\,730 + 19\,770 = 913\,500$، والمقام $923\,520 + 85\,680 =
1\,009\,200$، ومنه $Y = 0.91$. والأمومي عند الضغط نفسه: البسط
$893\,730 + 197\,730 = 1\,091\,460$، والمقام $923\,520 + 856\,830
= 1\,780\,350$، ومنه $Y = 0.61$. فعند الضغط نفسه يمسك الدم الجنيني أكسجين
أكثر، فيجري الأكسجين من الأم إلى الجنين عبر المشيمة.
\textbf{19.} $(10^{5}\,\text{nm}/6\,\text{nm})^{3} = 4.6\times 10^{12}$
خلية؛ و$1.9\times 10^{13}$ جزيء.
\textbf{20.} $d = 0.1/(2\sin 14.5^{\circ}) = \qty{0.20}{nm}$: نعم،
توضع كل سلسلة جانبية وكل ذرة.
\textbf{21.} $\sin\theta = 0.1/0.3$، ومنه $\theta = \qty{19.5}{\degree}$. وفي
البلورة غير المنتظمة لا تكون الجزيئات في اصطفاف دقيق، فلا تعود الأمواج
المستطارة عند الزوايا الكبيرة --- وهي التي تشفّر التفصيل الدقيق ---
تتجمع في الطور، فتخبو البقع عند الزوايا الكبيرة؛ ولا تنجو إلا
بقع \omterm{def:b3:structural-biology:methods}{التمييز} المنخفض.
\textbf{22.} تنصيف $d$ يحتاج ضعف الإشارة، ومن ثم أربعة أضعاف
الجسيمات: أي نحو \num{40000} (وأكثر عمليًا، إذ تدخل أخطاء أخرى).
\textbf{23.} البروتينات الغشائية، وهي نادرًا ما تتبلور من محاليل
المنظفات ويمكن تصويرها في قرص نانوي دهني؛ والمعقّدات الكبيرة
المرنة --- الريبوزومات، والجسيمات المشذِّبة، والقنوات الأيونية مع
منظماتها --- التي يمنع تغايرها التبلور لكن يمكن فرز
جسيماتها في أصناف انتظامية.
\textbf{24.} الهيم المرتبط والأكسجين وهندستهما الدقيقة، و
الماء والأيونات المنتظمة، والانتظام الثاني (T في مواجهة R) و
البرهان على صحة \omterm{def:b3:structural-biology:levels}{الطية} --- فالتنبؤ نموذج واحد بلا رابطة ولا فحص
تجريبي.
\textbf{25.} تفريغ $0.62$ من السعة (وبموقع واحد $0.36$)؛ وزمن
ليفينثال نحو $10^{52}$ سنة؛ وخريطة عند \qty{0.20}{nm}.
\end{solution}
